\documentclass[11pt]{article}
% Omit pdfTeX's path-dependent identifier for reproducible archive builds.
\pdftrailerid{}
\usepackage[T1]{fontenc}
\usepackage{lmodern}
\usepackage[margin=1in]{geometry}
\usepackage{amsmath,amssymb,amsthm,mathtools}
\usepackage{booktabs,array,microtype}
\usepackage[numbers,sort&compress]{natbib}
\usepackage{xurl}
\usepackage[colorlinks=true,linkcolor=blue,citecolor=blue,urlcolor=blue]{hyperref}
\hypersetup{pdftitle={Concentrating perishable preparation for deadline readiness},pdfauthor={Ruge Lin}}
\newtheorem{theorem}{Theorem}[section]
\newtheorem{lemma}[theorem]{Lemma}
\newtheorem{proposition}[theorem]{Proposition}
\newtheorem{corollary}[theorem]{Corollary}
\theoremstyle{remark}
\newtheorem{remark}[theorem]{Remark}
\newtheorem{example}[theorem]{Example}
\allowdisplaybreaks[1]
\title{Concentrating perishable preparation for deadline readiness}
\author{Ruge Lin}
\date{22 September 2026\\[4pt]\small Working research draft}

\begin{document}
\maketitle
\begin{abstract}
Preparation for a future handoff can deteriorate before it is used. We study
how to distribute such preparation across independent modules when a shared
source must be removable within a fixed deadline after any request. Module
$i$ has bounded preparation $p_i$, proportional deterioration $\gamma_i p_i$,
and releases a nonnegative source reservation after an instantaneous
independent transfer. We prove that every deadline-feasible state minimizing
total maintenance $\sum_i\gamma_i p_i$ has at most one partially prepared
module. The coefficients may differ, and some modules may initially be
inaccessible at the available preparation rate. A supporting serial reduction
covers arbitrary measurable, parallel, and preemptive allocations. The
concentration theorem yields an explicit finite frontier using
$O(n3^n)$ real scalar evaluations; unequal coefficients can require cold
modules before the partial module. Exact counterexamples delimit the result:
quadratic deterioration can require multiple partial modules, and finite
cold startup can sustain critical readiness without reaching its unique
minimum-maintenance state. A request-time precision analysis proves stability
of minimum upkeep below the available budget in the vanishing-deficit limit,
but exhibits a nonvanishing cost jump at a critical full endpoint.
An exchange using concave completion-time maps proves that concentration
itself survives every feasible positive preparation tolerance.
The results concern an ideal transfer interface;
the structural claims concern preparation inventories and their charged
maintenance resources, rather than a particular replication protocol.
\end{abstract}

\noindent\textbf{Development status.}
This is a working draft for mathematical and contribution review.
Its proofs have received internal adversarial checks and its exact examples
have small local verification scripts. External review, a validated service
implementation, and a settled novelty assessment are outstanding. No
submission or formal release is represented by this draft.

\section{Introduction}

A system that must relinquish control on request can spend resources in
advance to make that request easier to satisfy. If preparation loses value
while waiting, this creates a recurring cost even when the request never
arrives. The allocation question is concrete: should the system maintain
many partly prepared modules, or concentrate preparation into a few complete
modules? A deadline on the eventual exit couples that choice to the order
and timing of work performed after the request.

We formulate this question for independent modules under a deliberately
explicit interface. A module transfers when its preparation reaches a fixed
threshold; transfer is instantaneous, removes dependence on the source, and
releases a prescribed nonnegative source reservation. Receiver resources are
provided independently. Before a request, a fixed spare budget is shared by
preparation and optional work, and no transfer occurs. After a request,
optional work stops and released reservations can accelerate the remaining
preparation. Loss may follow any path within a specified state-dependent
envelope, including simultaneous maximum loss. The source cooperates during
the allowed exit interval: sudden source failure is a different problem.

The main result concerns positive proportional loss, with a possibly
different coefficient for each module. Every state minimizing the
maintenance needed to meet a given deadline has all but at most one of its
coordinates at an endpoint. This conclusion is stronger than existence of a
concentrated optimizer. It holds without assuming that the initial source
budget exceeds each module's full-maintenance rate. The proof varies two
partial coordinates while keeping an attaining exit schedule's later
completion time fixed. A strict concavity or strict monotonicity alternative
rules out both coordinates being interior. The argument uses the particular
maintenance prices supplied by deterioration; it does not extend without
proof to arbitrary investment prices.

This structure reduces a continuous initial-state optimization to a finite
description: a full subset, a cold prefix, at most one partial module, and a
cold suffix. The corresponding subset recurrences give the exact mathematical
frontier with an exponential real-evaluation bound. A cold prefix is
substantive when coefficients differ: an explicit example has a smaller
global maintenance optimum than any partial-first policy. The reference
floating implementation is an aid to inspection, not a certified oracle
for transcendental comparisons.

Several boundaries clarify what the theorem says. First, monotone smooth
deterioration still permits serial exit schedules but need not concentrate
the minimizing initial state. A rational quadratic-loss example separates a
two-partial feasible state from every concentrated state. Second, a minimum
maintenance value characterizes initialized long-run performance, but does
not by itself decide cold startup. At exact budget equality, a three-module
trajectory can remain ready indefinitely after finite warmup while converging
to a unique stationary optimum that no finite cold warmup can robustly reach or
dominate. Its finite reserve pays an integrable excess maintenance cost.

Third, exact threshold attainment can matter even when cold exit is finite.
We require readiness after a bounded preparation shortfall revealed at the
request and characterize the one-sided exit-time envelope. If minimum upkeep
is strictly below the normal budget, its limiting value is unchanged and
sufficiently small positive tolerances remain affordable. At a critical
endpoint, however, an exact two-module example raises limiting upkeep from
$1$ to $7/6$. This jump persists as tolerance tends to zero; it is stronger
than the ordinary positive cost of maintaining any safety margin.
The nominal optimizer nevertheless remains concentrated at every feasible
positive tolerance, even with unequal decay rates. A stronger exchange proves
this by allowing prepared modules between the two coordinates being varied.
Thus robustness of the optimizer's structure and continuity of its cost are
different properties.

These questions connect to established deteriorating-work scheduling,
controllable processing times, linear resource allocation, and storage
inequalities. In particular, common-rate fixed-order concentration has a
standard linear-allocation interpretation, and nonpreemption under
deterioration is an established research topic. Section~\ref{sec:related}
states the inspected correspondences. The contribution being assessed is the
unequal-rate minimum-maintenance structure and its precise limits in this
model, rather than a general new law of controller removal.

\section{Model and the readiness objective}
\label{sec:model}

Let \(N=\{1,\ldots,n\}\) be a finite collection of independent modules.
Module \(i\) has preparation capacity \(M_i>0\) and releases source capacity
\(a_i\geq0\) when it transfers. The initial spare source capacity is \(s\geq0\).
Preparation \(p_i\) belongs to \([0,M_i]\). Before a removal request, no
module transfers, and nonnegative preparation allocations \(v_i\) and
optional work \(u\) satisfy
\begin{equation}
 u+\sum_{i\in N}v_i\leq s.
 \label{eq:normal-budget}
\end{equation}
Following a request, optional work stops. If \(D\) is the transferred set,
the available preparation capacity is
\begin{equation}
 b(D)=s+\sum_{j\in D}a_j.
 \label{eq:released-capacity}
\end{equation}
A module can transfer when its preparation equals \(M_i\). Transfer is
instantaneous, removes that module's preparation obligation, and releases
\(a_i\) immediately. Releases and transfers have no cross-module constraints.

The main structural results concern proportional losses
\(g_i(p)=\gamma_i p\), with \(\gamma_i>0\). The scheduling reduction will be
proved for the larger class
\begin{equation}
 g_i:[0,M_i]\longrightarrow[0,\infty)
 \quad\text{nondecreasing and Lipschitz},\qquad g_i(0)=0.
 \label{eq:smooth-class}
\end{equation}
Allowed actual losses are measurable and obey
\(0\leq\rho_i(t)\leq g_i(p_i(t))\).
The uncertainty set includes every simultaneous maximum feedback
\(\rho_i=g_i(p_i)\); it is a product constraint on the actual states.
Under maximum feedback, before transfer the dynamics are
\begin{equation}
 \dot p_i=
 \begin{cases}
 v_i-g_i(p_i),&0\leq p_i<M_i,\\
 \min\{v_i-g_i(M_i),0\},&p_i=M_i.
 \end{cases}
 \label{eq:reflected-dynamics}
\end{equation}
There is no lower reflection term because \(g_i(0)=0\) and \(v_i\geq0\).
The corresponding upper-reflected dynamics apply under smaller losses.

Policies are deterministic and causal, with measurable allocations. They may
prepare modules in parallel, interrupt preparation, and adapt to observations.
Every guarantee below quantifies over all allowed loss histories. In
particular, maximum feedback is an admissible adversary for every policy,
rather than a fixed numerical loss trace shared by different trajectories.
Scalar comparison implies that larger initial preparation and smaller losses
cannot worsen a feasible exit strategy. Transfers can likewise be performed
earlier without harming another module: capacity is nondecreasing and the
modules are independent.

Write \(F(p)\) for the infimum, over such policies starting at \(p\), of their
worst-case time to transfer all modules. Infinite values are allowed. For a
finite deadline \(H\geq0\), define the ready set and its minimum upkeep by
\begin{equation}
 \mathcal R_H=\{p\in\textstyle\prod_i[0,M_i]:F(p)\leq H\},
 \qquad
 C(H)=\min_{p\in\mathcal R_H}\sum_i g_i(p_i).
 \label{eq:readiness-cost}
\end{equation}
Attainment of this minimum is established below. A uniformly ready normal
policy must belong to \(\mathcal R_H\) at every counterfactual request time,
after every admissible normal loss history.

The model presupposes a sufficient transfer state and separately paid
receiver resources. It does not derive a proportional loss envelope from a
network or replication protocol. Initial preparation is paid before the
uniform-readiness guarantee begins unless a separate warmup construction is
given. Positive handoff latency, delayed capacity release, and interdependent
transfers are outside these assumptions.

All normal preparation input is charged to \eqref{eq:normal-budget}.
Accordingly, $C(H)$ measures maintenance through this specified source
bottleneck, excluding receiver provisioning and independent continuing
operation. The prohibition on normal-mode transfer is an additional service
requirement, not a consequence of preserving essential service alone.
Setting every $a_i=0$ gives an exact alternative interpretation: independent
perishable jobs must remain uncompleted until a command permits absorbing
completion. States, controls, losses and deadlines are unchanged. Thus
positive resource release and controller removal are not necessary for the
concentration property; its minimal interpretation is preparation for a
future completion command.

\section{Serial exit from arbitrary preparation}
\label{sec:seriality}

A full-capacity serial schedule transfers initially full modules at time zero
and then assigns the entire available capacity to one remaining module until
it transfers. Its order is chosen as part of the optimization.

\begin{lemma}[Postponing scalar input]
\label{lemma:postponement}
Extend a function \(g\) satisfying \eqref{eq:smooth-class} constantly beyond
\(M\). Consider the uncapped equation \(y'=v-g(y)\) with the same initial
condition under nonnegative bounded measurable inputs \(v_{\rm old}\) and
\(v_{\rm new}\). If
\[
 D(t)=\int_0^t(v_{\rm new}-v_{\rm old})(r)\,dr\leq0
 \quad(0\leq t\leq T),\qquad D(T)=0,
\]
then \(y_{\rm new}(T)\geq y_{\rm old}(T)\).
\end{lemma}
\begin{proof}
Both solutions remain nonnegative. Set \(z=y_{\rm new}-y_{\rm old}\) and
\[
 \alpha(t)=
 \begin{cases}
 \dfrac{g(y_{\rm new}(t))-g(y_{\rm old}(t))}{z(t)},&z(t)\ne0,\\
 0,&z(t)=0.
 \end{cases}
\]
Monotonicity and Lipschitz continuity give
\(0\leq\alpha(t)\leq\operatorname{Lip}(g)\) almost everywhere. Thus
\(z'=v_{\rm new}-v_{\rm old}-\alpha z\), and variation of constants gives
\[
 z(T)=\int_0^T k(t)D'(t)\,dt,\qquad
 k(t)=\exp\left(-\int_t^T\alpha(r)\,dr\right).
\]
Since \(k'=\alpha k\geq0\), integration by parts yields
\(z(T)=-\int_0^T k'(t)D(t)\,dt\geq0\).
\end{proof}

The uncapped equation is only a comparison device. In a constructed physical
schedule, transfer a module at its first hit of \(M\) and omit its subsequent
planned allocations. This increases available capacity. A module that has
not hit \(M\) agrees with its uncapped comparison trajectory.

\begin{theorem}[Arbitrary-state seriality]
\label{thm:seriality}
Under \eqref{eq:smooth-class}, every finite feasible maximum-loss exit schedule
can be replaced by a full-capacity serial schedule whose final completion is
no later. Consequently, if \(F(p)<\infty\), an optimizing serial order attains
the robust value \(F(p)\). The optimum is attained among finitely many orders.
No preservation of an arbitrary schedule's completion order is asserted.
\end{theorem}
\begin{proof}
Initially full modules can transfer immediately. Let \(b\) be the resulting
capacity, and consider the remaining subfull modules. We may likewise transfer
every module at its first hit of its target: delaying a transfer cannot help,
because its release is nonnegative and its remaining loss is then removed.
If \(b=0\), none can
increase, so no finite exit is possible. Otherwise take a finite schedule
under maximum loss, and let \(i\) be a first module to complete, at time \(T\).
Capacity is \(b\) before that time.

A subfull module can complete at constant capacity \(b\) only if
\(b>g_i(M_i)\). Indeed, if \(b\leq g_i(M_i)\), a root of \(b-g_i\) is an
equilibrium barrier for the full-rate scalar equation. Lipschitz uniqueness
prevents finite-time arrival from below or crossing of that root. Starting
above a lower root gives nonpositive full-rate drift and cannot reach the
larger target. Allocating less than \(b\) cannot improve this comparison.
Conversely, \(b>g_i(M_i)\) gives a strictly positive drift throughout the
preparation interval.

The time to prepare \(i\) immediately at rate \(b\) is therefore
\[
 \tau=\int_{p_i(0)}^{M_i}\frac{dq}{b-g_i(q)}.
\]
Let \(V_i=\int_0^T v_i(t)\,dt\) in the original schedule. For
\(W_i(p)=\int_0^p dq/(b-g_i(q))\), its trajectory satisfies
\[
 \frac{d}{dt}W_i(p_i(t))
 =\frac{v_i-g_i(p_i)}{b-g_i(p_i)}
 \leq\frac{v_i}{b},
\]
because \(0\leq v_i\leq b\) and \(g_i\geq0\). Integration gives
\(V_i\geq b\tau\). If \(W\) denotes all work allocated to the other modules
before \(T\), the shared budget gives
\begin{equation}
 W\leq b(T-\tau).
 \label{eq:first-stage-work}
\end{equation}

We now postpone all those other modules' work to \([\tau,T]\), retaining
each module's total allocation and never increasing its cumulative allocation
before \(T\). For \(j\ne i\), let
\(E_j=\int_0^\tau v_j(t)\,dt\). On \([\tau,T]\), retain its old allocation
and define spare rate
\[
 h(t)=b-\sum_{j\ne i}v_j(t)\geq0.
\]
Equation \eqref{eq:first-stage-work} implies
\[
 \int_\tau^T h(t)\,dt
 =b(T-\tau)-W+\sum_{j\ne i}E_j
 \geq\sum_{j\ne i}E_j.
\]
Order the other modules arbitrarily and write
\(B(t)=\int_\tau^t h(r)\,dr\). Cumulative repayment to module \(j\) is
\[
 A_j(t)=\min\left\{E_j,
       \max\left\{0,B(t)-\sum_{\ell\ {\rm before}\ j}E_\ell\right\}\right\}.
\]
These absolutely continuous functions satisfy \(A_j(T)=E_j\) and
\(\sum_j A_j'\leq h\) almost everywhere. Give other modules zero allocation
before \(\tau\), and allocation \(v_j+A_j'\) afterward. Their total rate is at
most \(b\). For each module the new cumulative input is no larger than the
old one at every intermediate time and is equal at \(T\).

Prepare \(i\) at rate \(b\) until \(\tau\) and transfer it. Its release makes
the postponed allocations feasible. Lemma~\ref{lemma:postponement} shows that
at \(T\) every other module has at least its original preparation, unless it
has already completed, which is better. Stop such newly completed modules
and omit their later planned allocations. Simultaneous original completions
cause no difficulty: all their targets have been reached by \(T\).

At \(T\) the constructed state and completed set dominate the originals.
Replay the original suffix while omitting allocations to already completed
modules. Scalar comparison and nondecreasing capacity preserve domination.
The resulting schedule finishes no later than the original and has a
full-capacity first stage. Apply the same argument to its residual schedule
at \(\tau\). Induction on the number of remaining modules gives a serial
schedule no slower than the original.

For a fixed initial vector there are finitely many serial orders. Every
finite stage has a unique full-rate completion time, and inaccessible orders
are assigned infinity. The least finite order value is therefore attained.
Any robust policy must face simultaneous maximum feedback, so its
worst-case time is at least this value. Conversely, execute an optimizing
order against a virtual maximum-loss trajectory. Under smaller losses,
transfer each module no later than its virtual completion; greater
preparation and earlier capacity release preserve the comparison for the
remaining stages. This policy attains the same bound for every allowed loss
history and proves the robust assertion.
\end{proof}

\begin{lemma}[Compact readiness]
\label{lemma:compactness}
For every finite \(H\geq0\), \(\mathcal R_H\) is nonempty and compact.
Consequently the minimum in \eqref{eq:readiness-cost} is attained.
\end{lemma}
\begin{proof}
The all-full state exits at time zero. To prove closedness without assuming
continuity of hitting times at rate barriers, fix a permutation and
nonnegative stage durations \(\ell_1,\ldots,\ell_n\) with
\(\sum_r\ell_r\leq H\). During stage \(r\), allocate the full capacity of its
completed prefix to the selected module; unserved modules decay passively.
Require every selected module to be full at its stage endpoint. A full
module may voluntarily wait for its specified stage, and a zero-duration
stage is allowed when it is already full.

Passive flows and fixed-capacity upper-reflected scalar flows are continuous
in initial state and elapsed time. This follows from the Lipschitz drift
estimate and scalar comparison; upper reflection does not increase the
distance between ordered scalar states. Each endpoint constraint is
therefore closed in the compact product of the preparation box and the
duration simplex. Its projection onto the initial states is compact.
Taking the finite union over permutations remains compact.

Every projected witness is a feasible schedule, including deliberate
delays of already full modules. Conversely, Theorem~\ref{thm:seriality}
supplies such a witness for every state in \(\mathcal R_H\). Thus this
projection equals the ready set, including initially full and inaccessible-rate
boundaries. The upkeep objective is continuous, so it attains its minimum.
\end{proof}

A uniformly ready normal policy guarantees long-run optional rate \(r\) if,
on every allowed loss history,
\[
 \liminf_{T\to\infty}\frac1T\int_0^T u(t)\,dt\geq r.
\]
The optimal guaranteed rate is the supremum of these guarantees over
uniformly ready normal policies, allowing paid initialization.

\begin{theorem}[Exact initialized upkeep]
\label{thm:maintenance}
With paid initialization, the maximum guaranteed long-run optional rate under
uniform \(H\)-readiness is \(s-C(H)\) when \(C(H)\leq s\).
If \(C(H)>s\), indefinite uniform readiness is impossible even with no
optional work.
\end{theorem}
\begin{proof}
On maximum feedback put \(Q=\sum_i p_i\). For a uniformly ready policy,
upper reflection and \eqref{eq:normal-budget} imply, almost everywhere,
\[
 Q'\leq \sum_i v_i-\sum_i g_i(p_i)\leq s-u-C(H).
\]
Hence for every duration \(T\),
\begin{equation}
 \int_0^T u(t)\,dt
 \leq [s-C(H)]T+Q(0)-Q(T).
 \label{eq:storage-bound}
\end{equation}
The bounded storage term proves the asymptotic upper bound on this allowed
history. If \(C(H)>s\), the same inequality contradicts \(u\geq0\) for
sufficiently large \(T\).

For attainment, initialize at a minimizing \(p^*\), allocate
\(v_i=g_i(p_i^*)\), and use \(u=s-C(H)\). Under maximum loss the target is
stationary. Under every smaller loss history scalar comparison preserves
\(p\geq p^*\). The ready set is upward closed by the same comparison, so
every request meets \(H\). This constant optional allocation attains the
bound. It does not assert that the actual state is stationary under smaller
losses, or that \(p^*\) is reachable from cold with the normal budget.
\end{proof}

\section{Concentration under unequal proportional losses}
\label{sec:concentration}

For the remainder of the core results let \(g_i(p)=\gamma_i p\), with all
\(\gamma_i>0\), and write \(d_i=\gamma_iM_i\).
A coordinate is \emph{partial} when \(0<p_i<M_i\).
If a serial stage for module \(i\) begins at global time \(t_0\), has capacity
\(b>d_i\), and the module's initial preparation was \(p_i\), its completion
time \(t_1\) satisfies
\begin{equation}
 e^{\gamma_i t_1}
 =\frac{b e^{\gamma_i t_0}-\gamma_i p_i}{b-d_i}.
 \label{eq:proportional-stage}
\end{equation}
The waiting state is \(p_i e^{-\gamma_i t_0}\). A subfull module cannot
complete at capacity \(b\leq d_i\), even if its initial preparation is
arbitrarily close to full.

\begin{theorem}[Unequal-rate concentration]
\label{thm:concentration}
For every \(s\geq0\) and finite \(H\geq0\), every minimizer defining \(C(H)\)
has at most one partial coordinate. Thus a minimizing state consists of a
fully prepared set, at most one partial module, and otherwise cold modules.
\end{theorem}
\begin{proof}
Suppose a minimizing state had at least two partial coordinates. Choose an
attaining serial schedule, transferring all initially full modules at time
zero. Let \(j,k\) be consecutive partial modules in its remaining order.
Every intervening module is cold. Denote the start and end of stage \(j\)
by \(t_0,t_j\), the end of stage \(k\) by \(t_k\), and their capacities by
\(b_j,b_k\). Put
\[
 D_j=b_j-\gamma_jM_j>0,\qquad
 D_k=b_k-\gamma_kM_k>0.
\]
These inequalities follow from actual finite stage feasibility, without any
assumption that the original capacity exceeds every full loss.

The total duration \(L\geq0\) of the intervening cold stages is independent
of their absolute start time and of initial preparations of \(j,k\).
Replace \(p_j\) by \(x\), keeping all other earlier stages fixed. Equation
\eqref{eq:proportional-stage} gives
\[
 t_j(x)=\frac1{\gamma_j}
 \log\frac{b_j e^{\gamma_jt_0}-\gamma_jx}{D_j}.
\]
Choose the replacement \(y(x)\) for \(p_k\) to preserve its original
completion time:
\begin{equation}
 \gamma_k y(x)=
 b_k e^{\gamma_k L}
 \left(\frac{b_j e^{\gamma_jt_0}-\gamma_jx}{D_j}\right)^{\gamma_k/\gamma_j}
 -D_k e^{\gamma_k t_k}.
 \label{eq:pair-compensation}
\end{equation}
At the original point this gives \(y=p_k\). Both coordinates are interior,
so a two-sided open interval of small variations remains inside their
physical bounds. Their stage durations remain positive by continuity
(or directly by \eqref{eq:proportional-stage}). All intervening cold stages
retain their durations. Once stage \(k\) ends, global time and the completed
set agree with the original schedule; every later unserved module has its
unchanged passive state at that same time. The suffix therefore reproduces
the original final deadline.

Only the selected pair's upkeep varies:
\begin{equation}
 c(x)=\gamma_jx+
 b_k e^{\gamma_k L}
 \left(\frac{b_j e^{\gamma_jt_0}-\gamma_jx}{D_j}\right)^{\gamma_k/\gamma_j}
 -D_k e^{\gamma_k t_k}.
 \label{eq:pair-cost}
\end{equation}
If \(\gamma_k<\gamma_j\), this is strictly concave on the open variation
interval. At least one sufficiently small one-sided variation strictly
decreases its value. If \(\gamma_k\geq\gamma_j\), differentiation gives
\begin{equation}
 c'(x)=\gamma_j-
 \frac{b_k\gamma_k}{D_j}
 e^{\gamma_k L}e^{(\gamma_k-\gamma_j)t_j(x)}<0.
 \label{eq:pair-derivative}
\end{equation}
Here \(b_k\geq b_j>D_j\), the times are nonnegative, and
\(\gamma_k\geq\gamma_j>0\). A small increase of \(x\) again strictly lowers
upkeep while \eqref{eq:pair-compensation} keeps \(y\) interior.
Either case contradicts minimality. This proves the assertion for every
minimum, without an iterative exchange or an unproved endpoint limit.
\end{proof}

\begin{corollary}[A restriction on an optimizer's cold prefix]
\label{cor:prefix}
Suppose a minimizing state has a partial module \(k\). In any serial order
for that state that meets deadline \(H\) and transfers its initially full
modules at time zero, every cold module \(j\) before \(k\) satisfies
\(\gamma_j>\gamma_k\).
\end{corollary}
\begin{proof}
All modules between \(j\) and \(k\) are cold by
Theorem~\ref{thm:concentration}. Apply the same variation at \(x=p_j=0\)
and at the interior coordinate \(p_k\). If \(\gamma_j\leq\gamma_k\),
\eqref{eq:pair-derivative} is strictly negative. A small positive \(x\)
keeps \(y(x)\) interior and strictly reduces cost while preserving the
deadline. This contradicts minimality.
\end{proof}

This is a restriction on schedules witnessing an upkeep minimum, not an
ordering rule for arbitrary partially prepared states. In particular, equal
coefficients force an optimizing partial coordinate to precede all cold
modules; unequal coefficients do not. Strict positivity is also substantive:
a zero coefficient would permit cost-free partial coordinates and invalidate
the assertion about every minimum.

\section{An explicit finite readiness frontier}
\label{sec:frontier}

We now eliminate continuous optimization over the initial preparation vector.
All statements in this section allow \(s=0\), zero releases, and inaccessible
initial stages. For a full set \(S\), write \(d(S)=\sum_{j\in S}d_j\).
Define the cold-stage duration
\begin{equation}
 c_i(b)=
 \begin{cases}
 \gamma_i^{-1}\log\!\bigl(b/(b-d_i)\bigr),&b>d_i,\\
 +\infty,&b\leq d_i.
 \end{cases}
 \label{eq:cold-stage}
\end{equation}
Let \(G(S)\) be the exact exit time after immediately transferring full
modules \(S\), with all others cold. The serial theorem gives
\begin{equation}
 G(N)=0,\qquad
 G(S)=\min_{\substack{i\notin S\\b(S)>d_i}}
 \{c_i(b(S))+G(S\cup\{i\})\}.
 \label{eq:cold-tail-dp}
\end{equation}
An empty minimum is infinity. An initially full module transfers without
using \eqref{eq:cold-stage}, even if its putative preparation rate would be
inaccessible.

Fix the initially full set \(S\). For \(T\subseteq N\setminus S\), let
\(L_S(T)\) be the earliest completion time for a cold prefix consisting
precisely of \(T\). Its recurrence is
\begin{equation}
 \begin{split}
 L_S(\varnothing)&=0,\\
 L_S(T)&=\min_{\substack{j\in T\\b(S\cup(T\setminus\{j\}))>d_j}}
 \left\{L_S(T\setminus\{j\})
       +c_j\bigl(b(S\cup(T\setminus\{j\}))\bigr)\right\}.
 \end{split}
 \label{eq:cold-prefix-dp}
\end{equation}
Unreachable prefix states have value infinity. A waiting partial module does
not affect these durations.

Choose disjoint \(S,T\) and \(i\notin S\cup T\). Set
\[
 t=L_S(T),\qquad b=b(S\cup T),\qquad R=G(S\cup T\cup\{i\}).
\]
Consider the candidate only if \(t,R<\infty\) and \(b>d_i\).
If module \(i\) has initial preparation \(q\), its completion after that
prefix occurs at
\begin{equation}
 \Phi_i(t;q,b)=\frac1{\gamma_i}
 \log\frac{b e^{\gamma_i t}-\gamma_i q}{b-d_i}.
 \label{eq:partial-completion}
\end{equation}
For \(0\leq q\leq M_i\), this time is at least \(t\), and
\begin{equation}
 \frac{\partial\Phi_i}{\partial t}
 =\frac{b e^{\gamma_i t}}{b e^{\gamma_i t}-\gamma_i q}>0.
 \label{eq:prefix-monotonicity}
\end{equation}
Thus the earliest prefix is optimal even though the waiting partial
preparation decays. The remaining cold suffix is autonomous, so its best
duration is \(R\).

Writing \([z]_+=\max\{0,z\}\), the least initial preparation for this
displayed strategy to meet \(H\) is
\begin{equation}
 q_{S,T,i}(H)=
 \frac{\left[b e^{\gamma_i t}
       -(b-d_i)e^{\gamma_i(H-R)}\right]_+}{\gamma_i}.
 \label{eq:frontier-candidate}
\end{equation}
Discard candidates with \(q_{S,T,i}(H)>M_i\); the others have cost
\(K_{S,T,i}(H)=d(S)+\gamma_i q_{S,T,i}(H)\).
The endpoint \(q=0\) is cold. The endpoint \(q=M_i\), when the formula is
defined, is a feasible schedule that voluntarily delays a full module.
A full state at \(b\leq d_i\) is represented by adding \(i\) to the initial
full set, never by interpreting a singular partial-stage formula.

\begin{theorem}[Unequal-rate frontier]
\label{thm:frontier}
For every finite \(H\geq0\),
\begin{equation}
 C(H)=\min\left(
 \{d(S):G(S)\leq H\}
 \ \cup\
 \{K_{S,T,i}(H):(S,T,i)\text{ is an eligible candidate}\}
 \right).
 \label{eq:explicit-frontier}
\end{equation}
Eligibility is exactly the finite-prefix, finite-suffix, strict-stage, and
preparation-box conditions above. The formula supplies an attaining initial
state and request schedule. It can be evaluated with
\(O(n3^n)\) real scalar evaluations and comparisons and \(O(2^n)\) working
storage, besides output certificates.
\end{theorem}
\begin{proof}
Every candidate is feasible by its displayed serial construction. Pure
full/cold candidates are feasible by \eqref{eq:cold-tail-dp}; the all-full
set is always one of them. Their costs therefore cannot be below \(C(H)\).

Conversely choose an attained minimum. By
Theorem~\ref{thm:concentration}, it has a full set \(S\) and at most one
partial coordinate. A state without a partial coordinate is already in the
first family. Otherwise let \(i\) be its partial coordinate and choose an
attaining serial order after immediately transferring \(S\). Its other
modules form a cold prefix \(T\) and a cold suffix. Every actual stage is
accessible. Replace the prefix by \eqref{eq:cold-prefix-dp} and the suffix
by \eqref{eq:cold-tail-dp}. Equation \eqref{eq:prefix-monotonicity} shows that
the earlier prefix cannot hurt; the shorter cold suffix cannot hurt either.
Finally reduce initial preparation of \(i\) to
\eqref{eq:frontier-candidate} if necessary. This gives a listed feasible
candidate with no greater cost. Equality follows.

The cold-tail table has \(O(n2^n)\) transitions. Across all full sets \(S\),
there are \(3^n\) disjoint pairs \((S,T)\). The total number of prefix
transitions is \(n3^{n-1}\), and the number of distinguished partial positions
has the same bound. Evaluate one full set's prefix table at a time and retain
only the best candidate, together with the global cold-tail table.
This uses \(O(2^n)\) working storage; an attaining order can be reconstructed
by storing or recomputing the selected prefix certificate.
\end{proof}

The bound counts evaluations and comparisons of real logarithmic and
exponential expressions. It is not a polynomial-time result, a rational-field
bound, or a certified bit-complexity claim. Floating implementations need
explicit comparison tolerances or interval certificates at eligibility
boundaries. Corollary~\ref{cor:prefix} permits pruning prefixes containing
\(\gamma_j\leq\gamma_i\), but the unpruned formula is already exact.
Equation \eqref{eq:explicit-frontier} also gives
\(C(0)=\sum_i d_i\), and \(C(H)=0\) exactly when \(G(\varnothing)\leq H\).

\begin{corollary}[Common coefficient]
\label{cor:common-rate}
If all coefficients equal \(\gamma>0\), only empty cold prefixes are needed
in \eqref{eq:explicit-frontier}. The evaluation count reduces to
\(O(n2^n)\). For rational model parameters and rational
\(X=e^{\gamma H}\), all operations can be performed exactly in the rational
field.
\end{corollary}
\begin{proof}
Corollary~\ref{cor:prefix} excludes a cold predecessor of an optimizing
partial module. Put \(P(S)=e^{\gamma G(S)}\), with infinite values retained.
Then \(P(N)=1\), and accessible transitions obey
\[
 P(S)=\min_{\substack{i\notin S\\b(S)>d_i}}
       \frac{b(S)}{b(S)-d_i}P(S\cup\{i\}).
\]
For an empty-prefix candidate the required preparation becomes
\[
 \frac1\gamma
 \left[b(S)-(b(S)-d_i)\frac{X}{P(S\cup\{i\})}\right]_+.
\]
Keep the same accessibility and box tests and the full/cold candidates
\(P(S)\leq X\). These are rational expressions and comparisons under the
stated arithmetic assumptions. Rational \(H\) alone does not make \(X\)
rational.
\end{proof}

\begin{proposition}[A cold prefix can improve the global optimum]
\label{prop:cold-prefix}
With unequal coefficients, restricting an optimizing partial module to the
first nonfull position can strictly increase the minimum upkeep, already for
two modules with every initial cold stage accessible.
\end{proposition}
\begin{proof}
Take \(s=1\), \(H=\log3\), and
\[
 (M_A,a_A,\gamma_A)=(1/2,3/10,1),\qquad
 (M_B,a_B,\gamma_B)=(1/2,1/10,1/2).
\]
For a state with only the indicated coordinate partial, the four
order/coordinate choices require upkeep
\[
\begin{array}{c|c|c}
 \text{partial coordinate}&\text{order}&\text{required upkeep}\\ \hline
 A&AB&29/1352\\
 B&AB&(26\sqrt2-21\sqrt3)/20\\
 B&BA&1-\tfrac34\sqrt{18/11}\\
 A&BA&7/45.
\end{array}
\]
These follow directly from \eqref{eq:frontier-candidate}. For example,
cold A takes \(\log2\), after which B has capacity \(13/10\);
cold B takes \(2\log(4/3)\), after which A has capacity \(11/10\).
The second cost is positive and rationalizes to
\[
 C_*=\frac{29}{20(26\sqrt2+21\sqrt3)}<\frac{29}{1352}.
\]
For the strict inequality use \(\sqrt2>7/5\), \(\sqrt3>17/10\).
Also \(\sqrt{18/11}<13/10\), so the third cost exceeds
\(1/40>29/1352\); the fourth is larger as well. Every initially full
module costs at least \(1/4\). The all-cold orders miss the deadline because
their exponentiated durations are \(1352/441>3\) and \(88/27>3\).

Concentration and seriality make these comparisons exhaustive. Hence the
global minimum is \(C_*\), attained with
\(p_A=0\), \(p_B=(26\sqrt2-21\sqrt3)/10\), and cold A before partial B.
The best partial-first candidate costs \(29/1352>C_*\).
\end{proof}

\section{Boundaries of concentration and initialization}
\label{sec:boundaries}

The concentration theorem identifies the shape of a minimum-cost ready
state. Two separations delimit that statement. Nonlinear loss can make
distributed initial preparation strictly cheaper than every concentrated
state. Even under proportional loss, entering a sustainable ready regime
need not entail reaching any minimum-cost state in finite time.

\subsection{Smooth nonlinear loss can defeat concentration}
\label{sec:nonlinear-boundary}

Replace the proportional envelope by
$0\leq\rho_i(t)\leq g_i(p_i(t))$, keeping the other service, capacity,
and transfer assumptions unchanged. The next example uses identical
nonnegative, nondecreasing Lipschitz functions with $g_i(0)=0$.

\begin{proposition}[Nonlinear separation]
\label{prop:quadratic-failure}
For two identical modules with
\[
 M_1=M_2=1,\qquad a_1=a_2=\frac1{100},\qquad s=10,
 \qquad g_i(p)=p^2,\qquad H=\frac{23}{500},
\]
the ready state $(4/5,4/5)$ has maintenance cost $32/25$, while every
concentrated ready state has strictly larger cost.
\end{proposition}

\begin{proof}
Prepare module 1 at rate ten. Its loss is at most one, so its preparation
time is at most $1/45$. During that interval module 2 loses at most
$(16/25)/45$ and retains at least $884/1125$. After the first transfer,
capacity is $1001/100$ and the second module's net growth is at least
$901/100$. Its remaining deficit is at most $241/1125$. Hence the serial
policy's total time is at most
\[
 \frac1{45}+\frac{964}{40545}
 =\frac{373}{8109}<\frac{23}{500},
 \qquad
 \frac{23}{500}-\frac{373}{8109}=\frac7{4054500}>0.
\]
These are robust upper bounds, not exact quadratic traversal times.
Maintenance allocations $16/25$ to each coordinate preserve this target
under every allowed loss history.

For the converse, a concentrated state with no full module has total
initial preparation at most one. Transfers preserve the sum of completed
work and unfinished preparation, and nonnegative loss cannot increase it.
At least one additional unit of useful work is therefore necessary. Before
the last transfer, capacity is at most $1001/100$, whereas the deadline
permits at most
\[
 \frac{1001}{100}\frac{23}{500}
 =\frac{23023}{50000}<1
\]
units of allocated work. Such a state cannot be ready, under any schedule.

If one module is full and the other has preparation $r$, the same work
bound requires
\[
 r\geq 1-\frac{23023}{50000}=\frac{26977}{50000}.
\]
Its maintenance consequently satisfies
\[
 1+r^2-\frac{32}{25}
 \geq \frac{27758529}{2500000000}>0.
\]
Both modules full cost two. These cases exhaust concentrated states and
prove the claim.
\end{proof}

Thus smoothness, monotonicity, positive releases, and ample initial
capacity do not suffice for concentration. The proportional assumption is
substantive; equality of the proportional coefficients is not. The example
does not compute the nonlinear minimum itself.

\subsection{Critical readiness converges without necessarily becoming stationary}
\label{sec:startup-boundary}

We return to positive proportional coefficients. A normal warmup starts
from $p=0$, uses the same spare capacity $s$, and makes no transfers. A
removal deadline need only hold after the warmup. Write
$L(p)=\sum_i\gamma_i p_i$ and $Q(p)=\sum_i p_i$.

\begin{proposition}[Critical convergence]
\label{prop:critical-convergence}
Suppose $C(H)=s>0$. On the simultaneous maximal-loss history, every normal
trajectory that remains ready after a finite time $T_0$ converges to one
minimum-cost ready state $p^*$. Moreover,
\begin{equation}
 \int_{T_0}^{\infty}\bigl(L(p(t))-s+u(t)\bigr)\,dt
 \leq Q(T_0)-\lim_{t\to\infty}Q(t)<\infty.
 \label{eq:critical-dissipation}
\end{equation}
Finite-time arrival at $p^*$ is not asserted.
\end{proposition}

\begin{proof}
Reflection and the capacity constraint give
$Q'\leq s-u-L(p)$. Readiness implies $L(p)\geq C(H)=s$, so $Q$ is
nonincreasing and bounded below; integration proves
\eqref{eq:critical-dissipation}. Each coordinate is uniformly Lipschitz,
with derivative bounded in absolute value by
$\max\{s,\gamma_iM_i\}$. The nonnegative function $L(p(t))-s$ is therefore
uniformly Lipschitz and integrable, which forces it to tend to zero. A
positive lower bound along infinitely many times would otherwise produce
disjoint intervals of fixed positive area.

Compactness of the ready set implies that the distance to its minimizing
set tends to zero. That set is finite by Theorem~\ref{thm:concentration}:
the full set and partial identity leave the partial amount uniquely
determined by $L(p)=s$. A continuous trajectory eventually close to a
finite set cannot switch between disjoint neighborhoods of its members.
It therefore converges to one member.
\end{proof}

Convergence is not a reduction of startup to reaching a stationary
minimum. The following exact construction distinguishes those tasks.

\begin{theorem}[Startup without a reachable minimizing state]
\label{thm:startup-separation}
There is a three-module instance with positive sizes, coefficients,
releases, and spare capacity, and with finite cold exit, such that
$C(H)=s>0$ and its unique minimizing state cannot be robustly reached or
dominated after any finite normal warmup from cold. Nevertheless a finite normal
warmup enters a policy that is robustly $H$-ready at every later time.
\end{theorem}

\begin{proof}
Take $s=1$ and
\[
\begin{array}{c|ccc}
 \text{module}&M_i&\gamma_i&a_i\\ \hline
 A&1/2&1&2\\
 B&1&2&33\\
 C&3&1/2&1
\end{array}
\qquad
 p^*=(1/2,1/4,0),\qquad
 H=\tfrac12\log(5/2)+2\log(24/23).
\]
The target costs one. Transferring A immediately supplies rate three;
B then takes $\tfrac12\log(5/2)$. Its release supplies rate 36, so cold C
takes $2\log(24/23)$. Hence $p^*$ is ready.

We first certify that it is the unique minimum. The exact comparison
\begin{equation}
 H<\tfrac12\log3<\log2
 \label{eq:startup-deadline-comparison}
\end{equation}
follows from $5\cdot24^4=1658880<1679046=6\cdot23^4$.
Every minimizing state costs at most one and is concentrated by
Theorem~\ref{thm:concentration}. Neither B nor C can be full at that cost.
The remaining cases are exhaustive:
\begin{enumerate}
\item If A is subfull, it must complete first: subfull B and C cannot
complete at initial rate one, below their full loss rates two and $3/2$.
Cold A itself takes $\log2>H$. If A is partial, concentration makes B and C
cold; after A, their possible next stages take respectively
$\tfrac12\log3$ and $2\log2$, each exceeding $H$.
\item If A is full and B is the possible partial coordinate, cold C first
takes $2\log2>H$. B must precede C, giving total time
$\tfrac12\log(3-2p_B)+2\log(24/23)$. Readiness requires $p_B\geq1/4$,
while the cost bound $1/2+2p_B\leq1$ requires $p_B\leq1/4$.
\item If A is full and C is the possible partial coordinate, cost at most
one requires $p_C\leq1$. Cold B first already exceeds $H$. C first takes
at least $2\log(5/3)>\tfrac12\log3>H$, using $625>243$.
\end{enumerate}
Thus $C(H)=1$ and $p^*$ is its sole minimizer. The pure full/cold cases are
included by zero partial preparation. Seriality makes this an all-policy
certificate. Cold exit is finite in order A, B, C, with time
$\log2+\tfrac12\log3+2\log(24/23)$.

Next, no finite normal warmup can reach or dominate $p^*$ robustly.
On maximal loss, define $Z=p_A+p_B-3/4$. Every normal policy satisfies
\[
 Z'\leq1-p_A-2p_B=-2Z+(p_A-1/2)\leq-2Z.
\]
Starting cold gives $Z(t)\leq-(3/4)e^{-2t}<0$. Domination of $p^*$ would
instead require $Z\geq0$. Allocating to C cannot evade this support-budget
bound.

Finally, construct a different normal policy. Allocate $(0,0,1)$ for
$2\log2$, reaching the virtual maximal-loss state $(0,0,1)$. Thereafter
allocate $(1/2,1/2,0)$ forever. At elapsed time $t$ after the first phase,
put $x=e^{-t/2}$. The virtual state is
\begin{equation}
 p_A=\tfrac12(1-x^2),\qquad
 p_B=\tfrac14(1-x^4),\qquad p_C=x.
 \label{eq:startup-trajectory}
\end{equation}
For a request from this state, use order A, B, C. Its first two completion
times obey
\[
 e^{t_A}=1+x^2,\qquad
 e^{2t_{AB}}=Y(x):=\tfrac52+6x^2+\tfrac72x^4.
\]
The final completion time $T$ satisfies
\[
 e^{T/2}=\frac{72Y(x)^{1/4}-x}{69},\qquad
 e^{H/2}=\frac{72(5/2)^{1/4}}{69}.
\]
Concavity of the fourth root yields, for $0<x\leq1/128$,
\begin{align*}
 72\bigl(Y(x)^{1/4}-(5/2)^{1/4}\bigr)
 &\leq \frac{18}{(5/2)^{3/4}}\left(6x^2+\tfrac72x^4\right)\\
 &<108x^2+63x^4\\
 &\leq x\left(\frac{108}{128}+\frac{63}{128^3}\right)<x.
\end{align*}
Consequently $T<H$ for every finite request time after a total warmup of
$2\log2+14\log2=16\log2$. Every allocation uses the original normal
budget, and no module transfers before a request.

Under smaller allowed losses, scalar comparison keeps actual preparation
at least the virtual state. Readiness is upward closed, and the request
schedule can only improve. Thus the construction gives the claimed robust
guarantee, not only a maximal-loss trajectory calculation.
\end{proof}

Along \eqref{eq:startup-trajectory}, after readiness begins,
\[
 L(p)=1+\tfrac12x-\tfrac12x^2-\tfrac12x^4>1,\qquad
 Q(p)=\tfrac34+x-\tfrac12x^2-\tfrac14x^4\downarrow\tfrac34.
\]
All three coordinates remain partial at every finite time. The positive
preparation in C decays more slowly than the deficits in A and B, and its
request-time benefit offsets those deficits. Its total excess maintenance
is finite, consistent with Proposition~\ref{prop:critical-convergence}.
Concentration classifies minimizing states; it does not prohibit these
nonstationary ready states of strictly higher instantaneous maintenance.

The separation preserves the initialization distinction in the model. A
paid minimum is sufficient for the recurring frontier, while a finite
normal warmup may enter readiness through a different trajectory. Neither
construction supplies the removal guarantee before its stated entry time.

\section{Preparation precision at full endpoints}
\label{sec:precision}

An exactly full coordinate can transfer immediately, whereas a strict
deficit may expose a rate barrier. We quantify this distinction for
$n\geq1$ modules with positive proportional coefficients. For a maintained
state $p$, suppose a request may reveal any state $z$ satisfying
$(p_i-\epsilon)_+\leq z_i\leq p_i$. Here $\epsilon>0$ is an absolute
preparation deficit in the common work units. This is a bounded perturbation
at the request, not a stream of additional normal-operation disturbances.
Monotonicity makes $p^\epsilon=((p_i-\epsilon)_+)_i$ its worst corner.
Define
\begin{equation}
 F_-(p)=\lim_{\epsilon\downarrow0}F(p^\epsilon),\quad
 C_-(H)=\inf_{F_-(p)\leq H}L(p),\quad
 C_\epsilon(H)=\inf_{F(p^\epsilon)\leq H}L(p),
 \label{eq:precision-costs}
\end{equation}
where $L(p)=\sum_i\gamma_i p_i$ and empty infima are infinite. The limit
exists by monotonicity. Maintenance is charged to the nominal state $p$.

\begin{proposition}[One-sided exit envelope]
\label{prop:precision-oracle}
Let $\mathcal P$ contain the permutations satisfying
$b(S_{k-1})>d_{\pi_k}$ at every stage, where
$S_{k-1}=\{\pi_1,\ldots,\pi_{k-1}\}$ and $d_i=\gamma_iM_i$.
For each such order, define $G_\pi(p)=t_n$ by
\begin{equation}
 t_0=0,\qquad
 t_k=\frac1{\gamma_{\pi_k}}
 \log\frac{b(S_{k-1})e^{\gamma_{\pi_k}t_{k-1}}
                  -\gamma_{\pi_k}p_{\pi_k}}
                 {b(S_{k-1})-d_{\pi_k}}.
 \label{eq:precision-recurrence}
\end{equation}
Then $F_-(p)=\min_{\pi\in\mathcal P}G_\pi(p)$.
If cold exit is finite, $F_-$ is finite and continuous on the closed
preparation box. Otherwise it is identically infinite.
Always $F\leq F_-$, with equality at every subfull state.
\end{proposition}

\begin{proof}
Every $p^\epsilon$ is subfull. Theorem~\ref{thm:seriality} and the
strict finite-time rate barrier therefore allow exactly the orders in
$\mathcal P$, independently of $\epsilon$. Equation~\eqref{eq:proportional-stage}
gives \eqref{eq:precision-recurrence}. Each recurrence is continuous on the
closed box, with positive denominators and nonnegative stage durations.
Consequently its finite minimum commutes with the limit. The same orders
are feasible from cold, so the family is nonempty exactly when cold exit
is finite. Monotonicity gives $F\leq F_-$, and the identical order formula
applies directly to $F$ at subfull states.
\end{proof}

The strict inequalities include initially full coordinates. They can have
zero-length stages at time zero when accessible; a full coordinate left
waiting decays and requires preparation again. If the full coordinates
admit a strictly accessible order from the empty set, placing them first
at zero time and following an optimal original schedule proves
$F_-(p)=F(p)$. In particular this holds when every full coordinate has
$d_i<s$.

\begin{proposition}[Maintenance below the normal budget]
\label{prop:precision-slack}
If $C(H)<s$, then $C_-(H)=C(H)$.
\end{proposition}

\begin{proof}
At an original minimum with upkeep below $s$, each full coordinate has
$d_i<s$. The preceding zero-length prefix construction makes that state
feasible for $C_-$, giving equality since $F\leq F_-$ implies $C\leq C_-$.
\end{proof}

At $C(H)=s$, failure of equality can occur only when every original
minimum is a full singleton of cost $d_i=s$. Any other minimum has all
full losses strictly below $s$ and supplies the same zero-length prefix.
The usual total-storage bound and maintenance construction apply to the
upward closed limiting ready set, giving optional rate $s-C_-$ when
$C_-\leq s$ and excluding indefinite readiness otherwise.

\begin{proposition}[Concentration with preparation tolerance]
\label{prop:precision-concentration}
Whenever the respective finite-deadline problem is feasible, every
minimizing state for $C_-(H)$ and every minimizing nominal state for
$C_\epsilon(H)$ has at most one partial coordinate. This holds for unequal
positive coefficients and at every source budget.
\end{proposition}

\begin{proof}
For the limiting problem use $q=p$ with caps $U_i=M_i$. For fixed positive
tolerance, an optimizer cannot have $0<p_i\leq\min(\epsilon,M_i)$:
replacing it by zero preserves its worst corner and lowers upkeep. Hence
use $q_i=(p_i-\epsilon)_+$ with caps $U_i=(M_i-\epsilon)_+$ and upkeep
$\sum_i\gamma_iq_i+\sum_{i:q_i>0}\gamma_i\epsilon$.
Zero q corresponds to zero p, while $q_i=U_i>0$ corresponds to $p_i=M_i$.
Both problems therefore use the strict-order oracle with capped q; the
additional positive-support charges stay constant under interior variations.
Their minima exist when feasible: the limiting feasible set is compact
by Proposition~\ref{prop:precision-oracle}, and the fixed-tolerance
nominal constraint is continuous by applying that oracle to its subfull
worst corner.

For fixed preparation q, a stage map $\phi(t)$ from
\eqref{eq:precision-recurrence} satisfies
\[
 \phi'\geq1,\qquad
 \phi''=-\gamma\phi'(\phi'-1)\leq0,\qquad \phi(0)\geq0.
\]
The composition G of any intervening stages is increasing and concave,
with $G(0)\geq0$. Its tangent inequality gives $G(t)\geq tG'(t)$.

Suppose two coordinates j,k are interior to their positive caps in an
attaining strict order. Keep the start of j fixed, vary its completion t,
and keep the completion of k fixed. All intervening preparations remain
fixed, with G mapping t to the start of k. Writing
$\alpha=\gamma_j$, $\beta=\gamma_k$, $D_j=b_j-d_j>0$, the pair upkeep
up to constant terms is
\[
 c(t)=-D_j e^{\alpha t}+b_k e^{\beta G(t)}.
\]
An open interval preserves both interior caps and positive stage durations;
the same completion of k reproduces the full suffix and its deadline.
At a stationary point,
$\alpha D_j e^{\alpha t}=\beta b_k G'e^{\beta G}$.
If $\beta G'\geq\alpha$, then $\beta G\geq\alpha t$, contradicting
this equality because $b_k\geq b_j>D_j$. Thus $\beta G'<\alpha$, and
\[
 c''(t)=\alpha D_j e^{\alpha t}
       \left(\beta G'+\frac{G''}{G'}-\alpha\right)<0.
\]
Every stationary point is a strict local maximum. The proposed minimum
cannot have two interior coordinates. Under the nominal conversion above,
this is exactly the asserted concentration.
\end{proof}

The worst-corner state can itself have several physically partial
coordinates: nominally full modules have preparation $M_i-\epsilon$ at
the request. The result concerns the nominal design, and does not remove
their post-request preparation or justify treating them as free transfers.

\begin{proposition}[Positive tolerance and its limit]
\label{prop:precision-limit}
If cold exit is finite and $H>0$, then
$\lim_{\epsilon\downarrow0}C_\epsilon(H)=C_-(H)$.
For every positive $\epsilon$, whenever $C(H)>0$,
\begin{equation}
 C_\epsilon(H)\geq C(H)+\gamma_{\min}\epsilon,
 \qquad \gamma_{\min}=\min_i\gamma_i>0.
 \label{eq:precision-margin}
\end{equation}
\end{proposition}

\begin{proof}
Monotonicity gives $F_-(p)\leq F(p^\epsilon)$, hence $C_\epsilon\geq C_-$.
Choose a minimizing state for $C_-$. If its time is strictly below $H$,
it is feasible at sufficiently small positive tolerance. Otherwise its
time equals $H>0$, so it is not all-full. Moving it an arbitrarily small
distance toward the all-full vector strictly lowers an attaining order's
time: increasing a coordinate strictly advances its own stage, and every
later recurrence is strictly increasing in its preceding time. The new
state has time slack and arbitrarily close upkeep. It is feasible for all
sufficiently small tolerances, proving the opposite limit inequality.

For \eqref{eq:precision-margin}, any feasible nominal state has
$q=p^\epsilon$ with $L(q)\geq C(H)>0$. Some $q_i>0$, so
$p_i=q_i+\epsilon$, while every other $p_j\geq q_j$. Thus
$L(p)\geq L(q)+\gamma_{\min}\epsilon$; take the infimum.
\end{proof}

The condition $H>0$ is essential: positive deficits make a zero deadline
infeasible even though the all-full state has $F_-=0$ when cold exit is
finite. For positive deadlines, strict optional slack survives sufficiently
small positive tolerance, although its exact optimal rate can decrease.
Every originally critical instance $C=s>0$ loses indefinite readiness at
any positive tolerance by \eqref{eq:precision-margin} and the storage
bound. That elementary margin obstruction is general; the next example's
distinctive feature is a nonvanishing limiting cost jump.

\begin{example}[A critical endpoint jump with finite cold exit]
\label{ex:precision-jump}
Take
\[
 s=\gamma_A=\gamma_B=1,\quad
 M_A=\tfrac12,\quad M_B=1,\quad a_A=a_B=1,
 \qquad H=\log(4/3).
\]
Only the order A,B is strictly accessible. Its recurrence is
$\exp F_-(p)=4-4p_A-p_B$, so cold exit takes $\log4$.
The original state $(0,1)$ transfers B immediately and finishes A at
capacity two in time $\log(4/3)$, whereas $F_-(0,1)=\log3$.
The limiting deadline requires $4p_A+p_B\geq8/3$, implying
$p_A+p_B\geq8/3-3p_A\geq7/6$. Equality holds only at
$(1/2,2/3)$. Every original ready state with $p_B<1$ obeys this same
constraint, while B-full states cost at least one. Consequently
\[
 C(H)=1=s,\qquad C_-(H)=\tfrac76>s.
\]
For fixed tolerance, both coordinates of $q=p^\epsilon$ must be positive
to satisfy $4q_A+q_B\geq8/3$. Nominal upkeep is then
$q_A+q_B+2\epsilon$, minimized by
$q_A=1/2-\epsilon$, $q_B=2/3+4\epsilon$. The latter obeys
$q_B\leq1-\epsilon$ exactly when $\epsilon\leq1/15$. Hence
\[
 C_\epsilon(H)=
 \begin{cases}
  7/6+5\epsilon,&0<\epsilon\leq1/15,\\
  +\infty,&\epsilon>1/15.
 \end{cases}
\]
For larger tolerance even the all-full nominal state is infeasible.
The limiting increase $1/6$ therefore persists as precision improves.
Two modules are necessary under finite cold exit: a single module requires
$s>d$ and has the continuous formula
$F(p)=\gamma^{-1}\log[(s-\gamma p)/(s-d)]$ through its full endpoint.
\end{example}

% Source inspection records: research/2026-09-22-pass11-prior-art.md and
% research/2026-09-22-pass15-assessment.md. Glazebrook 1992/1993 were inspected
% at abstract/metadata level only; no theorem-level separation is asserted.
\section{Related work and interpretation}
\label{sec:related}

\subsection{Deteriorating work and initial preparation}

Scheduling with deterioration already asks when delaying or interrupting a
job increases the work required to complete it. Glazebrook
\cite{Glazebrook1992,Glazebrook1993} studies sufficient conditions for
nonpreemptive or permutation policies in stochastic deterioration models.
A more directly inspectable repair-state formulation is given by Gehlot,
Sundaram, and Ukkusuri \cite{GehlotSundaramUkkusuri2020}. In their discrete
model, the selected component gains a fixed repair increment, other components
lose fixed health increments, and both failure and full repair are absorbing.
Their Theorem~1 removes preemptions when each deterioration increment is at
least the corresponding repair increment. Its proof postpones a partially
repaired component and compares the resulting service durations. Their
subsequent model incorporates precedence constraints and a finite deadline
\cite{GehlotSundaramUkkusuri2020Interdependent}.

Our supporting serial reduction uses the same broad idea that service moved
later suffers less deterioration. Its hypotheses differ: zero preparation is
recoverable, controls are continuous and divisible, loss depends on current
preparation, and completed modules may increase available capacity. We do not
present nonpreemption under deterioration as a new general principle. The
main optimization here concerns the \emph{initial preparation vector} and its
ongoing maintenance cost, after an exact exit-time characterization has been
established.

Controllable-processing-time scheduling is consequently another essential
comparison. Shioura, Shakhlevich, and Strusevich
\cite{ShiouraShakhlevichStrusevich2015} minimize linear compression cost under
a common deadline and develop algorithms for linear optimization over a
box-constrained submodular region. Their Section~3.1, Theorem~2, states the
standard greedy rule that contains our common-decay fixed-order optimization.
Indeed, if an order has deadline inequality
\[
 A_\pi-\sum_r w_r p_{\pi_r}\leq X,
 \qquad \sum_r w_rM_{\pi_r}=A_\pi-1,
\]
then the substitution $y_r=w_r(M_{\pi_r}-p_{\pi_r})$ gives
\[
 \max\sum_r\frac{\gamma}{w_r}y_r,
 \qquad 0\leq y_r\leq w_rM_{\pi_r},
 \qquad \sum_r y_r\leq X-1.
\]
Thus one-partial-coordinate existence in this specialization is ordinary
linear allocation. The strict ordering of the model's weights supplies the
additional preparation-prefix property. Taking a finite minimum over orders
does not change the attribution of the fixed-order optimization.

Wei, Wang, and Ji \cite{WeiWangJi2012} combine start-time deterioration with
controllable processing times of the form $a_j+bt-\theta u_j$. Their
Section~3 expands a prescribed order's completion time into a weighted affine
function of resource allocations. For identical module sizes $M$, common
decay $\gamma$, and zero releases, our serial recurrence becomes their
recurrence after setting $D=s-\gamma M>0$, $z=e^{\gamma t}-1$,
$b=\gamma M/D$, $a_j=\gamma M/D$, $\theta=\gamma/D$, and $u_j=p_j$.
This is an exact restricted timing correspondence; their formulation assumes
nonpreemption and optimizes weighted scheduling and resource costs.

Nonlinear resource dependence alone also does not distinguish the present
problem. Li, Wang, and Wang \cite{LiWangWang2015} consider durations
$((a_j/u_j)^k+bt)r^\alpha$ and derive optimal investment and sequencing from
separable convex resource terms and rearrangement. Our initial preparation
instead interacts with waiting through
\[
 \tau_i(t,p;b)=\frac{1}{\gamma_i}
 \log\frac{b-\gamma_i p e^{-\gamma_i t}}{b-\gamma_iM_i}.
\]
The mixed derivative with respect to $t$ and $p$ is strictly positive, whereas
their duration is additively separable in start time and investment. Directly
identifying their resource variable with initial preparation therefore does
not give this model.

A directly inspectable fixed-deadline nonlinear allocation appears in
Choi and Park \cite{ChoiPark2023}, Section~3, Lemma~3. Its proof minimizes
$\sum_j c_j u_j$ subject to $\sum_j(w_j/u_j)^k\leq\delta$, with positive
variables and coefficients. Convexity and the KKT conditions give
$u_j=(\lambda k w_j^k/c_j)^{1/(k+1)}$; the investments can all be interior.
Our fixed-order deadline geometry need not be convex. For example, with
$s=3$, zero releases, $M_1=M_2=1$, $(\gamma_1,\gamma_2)=(2,1)$ and
$H=\log2$, order $1,2$ requires
\[
 y\geq3\sqrt{3-2x}-4.
\]
The boundary points $(13/25,1/5)$ and $(3/8,1/2)$ are feasible, while their
midpoint $(179/400,7/20)$ is not: squaring the required positive-side
inequality would give $7578\leq7569$. An affine identification with that
convex investment subproblem is therefore impossible. This comparison
excludes that specified reduction; it does not exclude arbitrary coupled
transformations or establish priority over all investment scheduling results.

The unequal-rate concentration theorem is the structural question beyond
these fixed-order correspondences. Its deadline-preserving pair exchange
excludes two interior coordinates even without a common affine clock. The
cost coefficients are the decay coefficients themselves, rather than arbitrary
investment prices. Nonnegative capacity release preserves the exchange; the
result also includes zero releases. The quadratic-loss example identifies a
boundary of this particular concentration property, while leaving seriality
intact. It is not a claim that nonlinear resource models generally concentrate
investment.

\subsection{Maintenance, service, and initialization}

The passage from minimum maintenance to long-run optional throughput is a
storage argument. With $Q=\sum_i p_i$, a ready trajectory under maximum loss
satisfies $\dot Q\leq s-u-C(H)$. Integrating and using bounded storage gives
the average-throughput bound. This is the established connection between a
storage inequality and optimal steady operation discussed by Faulwasser
et al.\ \cite{FaulwasserEtAl2017}. We prove the particular reflected-system
inequality directly; the contribution under study is the readiness geometry
and its optimizer, not the averaging principle.

The model assumes an independent receiver and an instantaneous transfer
primitive. This is stronger than possession of a copied image. Practical
fault-tolerant execution must preserve externally acknowledged obligations:
Scales, Nelson, and Venkitachalam
\cite{ScalesNelsonVenkitachalam2010}, Sections~2.2--2.3, gate external output
on sufficient backup replay information and treat takeover separately.
Live migration likewise distinguishes pre-copy, source suspension, remaining
transfer, commitment, and activation \cite{ClarkEtAl2005}, Section~3.3.
These mechanisms demonstrate why bandwidth accounting alone does not imply
correct ownership transfer or a bounded service response.

Here proportional loss is a pathwise uncertainty envelope. Randomly located
record updates may motivate an expected proportional drift, but that
expectation does not establish the envelope for every update history. A
finite-delay implementation would also need an explicit contract for
acknowledgements, in-flight information, queuing, and fencing. The results
therefore describe perishable preparation under the stated ideal handoff
interface, rather than an implementation-independent migration bound.

Finally, maintaining an optimal state and reaching it from a cold system are
different questions. Paid initialization makes the stationary throughput
construction feasible whenever its maintenance fits the budget. Without that
allowance, strict maintenance slack permits finite warmup, while equality
requires a separate reachability analysis. The critical three-module
construction further separates finite entry into an indefinitely ready
trajectory from finite arrival at a stationary optimum: the former can be
possible when the unique optimum is unreachable. No readiness guarantee is
asserted during an unprepared warmup. This distinction is part of the service
contract, not an adjustment to the long-run objective.

\subsection{Why the charged resource and normal-mode restriction matter}

The source-only upkeep bound need not survive additional independent input.
For one module with $M=\gamma=1$, $s=2$, $a=0$ and $H=\log(3/2)$,
scalar comparison gives $F(p)=\log(2-p)$ and hence $C(H)=1/2$.
The declared model has optional optimum $3/2$. If an independently paid input
of rate $1/2$ can maintain the same preparation, initialize at $p=1/2$ and
use that input while assigning the entire source budget to optional work.
Every request still meets $H$ under the original rate-two exit strategy, and
normal optional throughput is two. Total maintenance has not vanished; its
cost has moved outside the source objective.

Likewise, if permanent transfers are allowed before the request, preserve
the same required service, and have no separately charged objective cost,
the initialized problem admits an immediate counterpolicy: initialize every
module fully and transfer them all. Every later request then has zero exit
time and needs no source preparation. A positive source-upkeep lower bound
therefore requires the stated normal-mode restriction or an objective that
accounts for the admissible alternatives. Neither restriction follows merely
from calling the eventual receiver independent. These counterpolicies delimit
the ideal service class; they do not assert that every physical receiver can
supply either additional control.


\section{Verification and remaining questions}

The analytical proofs carry the claims over continuous time and arbitrary
admissible policies. The accompanying repository preserves small
standard-library verification scripts for exact inequalities, finite
subset-versus-permutation comparisons, and explicitly labeled floating
identity checks. The quadratic and critical-startup certificates use rational
inequalities and closed-form expressions; no large simulation or empirical
fit enters their proofs. Verification counts are not independent experiments,
and successful runs do not establish model applicability or novelty.

Several questions remain separate from the proved structural result. The
closest older stochastic-deterioration results have not all been inspected
at theorem-and-proof level, so the present comparisons do not establish
literature-wide absence of an equivalent concentration theorem. A practical
service interpretation must justify the entire admissible interface, including
why normal operation cannot use receiver resources for earlier independent
maintenance, what preparation certifies, and when source reservations can
actually be released. A one-sided conservative loss envelope can justify a
safe policy without making the envelope's optimum a necessary physical cost.

The precision results cover a one-time request deficit, rather than recurring
normal-operation disturbances. Their concentration theorem does not supply
an unchanged frontier algorithm: nominally full modules now require processing
after the request. The general cold-entry viability problem and certified
numerical frontier evaluation are also not solved here.
These are concrete limitations of the present formulation and algorithms,
not conclusions supplied by finite verification.

\paragraph{Reproducibility.}
Editable sources, full research derivations, inspection scopes for cited
papers, failed conjectures, and verification commands are available in
\url{https://github.com/GoGoKo699/Safe-Disengagement-Limits}.
The accompanying \texttt{paper/README.md} identifies the build and review
procedure. Historical research checkpoints are preserved unchanged.

\bibliographystyle{plainnat}
\bibliography{references}
\end{document}
